% !TeX root = ../main.tex
\section{Putting three experiments on one axis}
\label{sec:method}

Comparing results across three implementations built years apart requires deciding what to
hold fixed. We fix three things: the grader, the pairing, and the test.

\subsection{One grader}
\label{sec:one-grader}
Generations~2 and~3 both used \primaryjudge{} (\texttt{gpt-4o-mini-2024-07-18}) at temperature
$0.1$ with a JSON-schema-constrained rubric. We verified that this is not merely nominal
agreement: constructing the Q1 and Q2 evaluation prompts from each generation's own
questionnaire module on the same transcript yields \textbf{byte-identical prompts, schemas and
label sets}. Generations~2 and~3 are therefore already on one measurement axis, and no
re-scoring was needed for either.

Generation~1 is the exception. It was graded by GPT-3.5 through a regex-parsed rubric, a
different model and a different extraction path. We therefore re-scored \textbf{all $1{,}344$
of its evaluation conversations} ($7$ iterations $\times$ $2$ look-ahead arms $\times$ $96$
personas, plus the $96$ base conversations) with the generation-3 oracle on Q1 and Q2:
$2{,}880$ calls, $0$ errors. Because the conversations already exist on disk, this required
no GPU time and no patient simulation --- only oracle calls --- and cost approximately
\$1.5.\footnote{Scores are written to a separate \texttt{\_crossgen} partition so that no
tracked generation-3 result can move as a side effect. The re-scoring tool is
\texttt{eda/tools/score\_crossgen.py}.}

This yields two independent readings of generation~1 --- its original GPT-3.5 scores and its
\primaryjudge{} re-grade --- which \S\ref{sec:regrade} uses to separate ``the grader changed''
from ``the experiment changed''. Unless stated otherwise, every number reported for every
generation from \S\ref{sec:regrade} onward is on the \primaryjudge{} axis.

\subsection{One pairing}
\label{sec:pairing}
All contrasts are \textbf{persona-paired}: the $\K{=}0$ and $\K{=}5$ conversations being
compared are with the \emph{same} simulated patient. Persona sampling is the dominant source
of variance in these evaluations (a persona's difficulty moves its score far more than the
therapist's iteration does), so pairing removes most of the noise and is what makes effect
sizes interpretable.

How the pairing is recovered differs by generation, and getting this wrong silently corrupts
every effect size:

\begin{itemize}\itemsep2pt
  \item \textbf{Generations 1 and 2} wrote \texttt{conversation\_\{i\}.csv} with a
  \emph{fixed} persona order. We verified this directly --- \texttt{conversation\_0} is the
  same patient (same name, age, presenting problem and cooperation level) across both $\K$
  arms and across every iteration. The file index therefore \emph{is} the persona.
  \item \textbf{Generation 3} reshuffles the 96 personas every iteration with seed
  $\text{seed}+k+1$. A file-index join would pair unrelated conversations. Its pairing is
  recovered by replaying the shuffle, through the existing analysis package rather than
  re-implemented here.
\end{itemize}

\subsection{One test}
\label{sec:test}
For each (generation, training oracle, optimizer, iteration, metric) cell we compute the
paired difference over the 96 personas,
\[
  \dlt \;=\; \bar{x}_{\K=0} - \bar{x}_{\K=5},
\]
reported with Cohen's $\dz$ (mean difference over the standard deviation of the differences)
and a two-sided Wilcoxon signed-rank $p$, Holm-corrected within each
(generation, oracle/optimizer, iteration) family across metrics.

\paragraph{Sign convention.}
\textbf{$\dlt$ is $\K{=}0$ minus $\K{=}5$, so a \emph{negative} $\dlt$ means look-ahead
helped.} This matches the convention of the generation-3 artifacts. It is the opposite
orientation to the one the ICLR paper used, so all conversions are done in the analysis code
and never in prose.

\paragraph{Matched iterations, not best-versus-best.}
Every contrast compares iteration $t$ of the $\K{=}0$ arm with iteration $t$ of the $\K{=}5$
arm. This is the substantive methodological difference from the published analysis, which
compared each arm's \emph{best} model. Selecting each arm's maximum before comparing them
inflates whichever arm is noisier or was run longer, and discards the trajectory. Matched
iterations also make the cost comparison meaningful: at equal $t$ the two arms have had the
same number of policy updates, and $\K{=}5$ has paid several times more per update to get
there.

\subsection{Reproducibility}
All tables and figures in this chapter are produced by a single script,
\texttt{analysis/crossgen.py}, which reads the three generations' score files and writes
\texttt{tables/*.md} and \texttt{figures/*.png}. Nothing is hand-computed; the claims ledger in
\texttt{NUMBERS.md} maps each number in the text to the artifact it came from. The oracle's own
reproducibility was measured in generation~3 --- $\mathrm{ICC}(2,1)$ between $0.86$ and $0.99$
across instruments, with Q1 and Q2 at $0.96$--$0.99$ --- so a difference of $0.10$ on \QQ{} is
comfortably above grader noise, and differences below ${\sim}0.05$ are not interpretable.
